\section{Diskussion der Ergebnisse}
hallo das ist ein test 
Die Entwicklung von Pablo zeigt, dass ein kompakter Desk-Companion-Bot mit sprachbasierter Interaktion, visueller Rückmeldung und modular erweiterbarem Funktionsumfang als integriertes System umgesetzt werden kann. Im Vergleich zu rein softwarebasierten Assistenzlösungen wurde mit Pablo bewusst ein physisch präsentes System realisiert, das nicht nur funktionale Aufgaben übernimmt, sondern auch als sichtbarer Begleiter im direkten Nutzungskontext auftritt. Die physische Präsenz, die audiovisuelle Interaktion und die mechanische Ausrichtung des Systems tragen dazu bei, dass die Interaktion stärker als persönlicher Austausch wahrgenommen werden kann.

Die wakeword-basierte Aktivierung stellt dabei eine geeignete Lösung für eine niedrigschwellige und gleichzeitig gezielte Nutzung dar. Das System verbleibt zunächst im Bereitschaftszustand und wird erst nach einer bewussten Ansprache aktiv. In Verbindung mit der Richtungsbestimmung und der Ausrichtung des Displays entsteht ein Interaktionsverhalten, das die Ansprechbarkeit des Systems räumlich sichtbar macht und den Companion-Charakter zusätzlich unterstützt.

Die Kombination aus freier Sprachverarbeitung über ein Large Language Model und strukturierten Funktionsaufrufen erweitert die Einsatzmöglichkeiten des Systems deutlich. Während allgemeine Anfragen flexibel verarbeitet werden können, ermöglichen klar definierte Module die zuverlässige Umsetzung konkreter Alltagsfunktionen. Die implementierten Module für Wetter, Notizen, Kalender und Timer decken typische Anwendungsszenarien im Alltag ab und eignen sich gleichzeitig als prototypische Nachweise für die Erweiterbarkeit der Architektur.

Die modulare Softwarearchitektur erweist sich hierbei als zentrale Grundlage des Systems. Durch die Trennung von Aktivierung, Spracherkennung, Sprachverarbeitung, Sprachausgabe und Hardwaresteuerung konnten einzelne Systembestandteile unabhängig voneinander entwickelt und angepasst werden. Insbesondere die Plugin-Schnittstelle stellt ein wesentliches Ergebnis dar, da sie eine funktionale Erweiterung des Systems ermöglicht, ohne den Kernprozess grundlegend verändern zu müssen. Dies ist besonders relevant für zukünftige Anpassungen an unterschiedliche Zielgruppen und Nutzungskontexte.

Auch das Hardwarekonzept unterstützt diese Zielsetzung. Die Auswahl kompakter und verbreiteter Komponenten wie Raspberry Pi, ReSpeaker, Waveshare-Display sowie separater Mikrocontroller für Servo- und Verstärkeransteuerung ermöglichte einen integrierbaren und gleichzeitig flexibel entwickelbaren Aufbau. Die Verteilung einzelner Aufgaben auf mehrere Recheneinheiten reduzierte die Belastung der Hauptrecheneinheit und erleichterte die Einbindung mechanischer und audiovisueller Funktionen.

Gleichzeitig zeigt das entwickelte System auch typische Grenzen eines prototypischen Produktes. Die Systemleistung hängt von der zuverlässigen Abstimmung mehrerer Hard- und Softwarekomponenten ab. Insbesondere die Kopplung von Audioverarbeitung, Sprachmodell, visueller Darstellung und mechanischer Bewegung erfordert eine robuste Synchronisation. Darüber hinaus ist die wahrgenommene Natürlichkeit der Interaktion nicht allein vom Sprachmodell abhängig, sondern ebenso von Reaktionszeit, Audioqualität, Wakeword-Zuverlässigkeit und konsistentem visuellen Feedback.

Insgesamt verdeutlichen die Ergebnisse, dass Pablo als modularer Desk-Companion-Bot erfolgreich als funktionaler Prototyp umgesetzt wurde. Besonders hervorzuheben sind die Verbindung aus physischer Präsenz und Sprachinteraktion, die pluginbasierte Erweiterbarkeit sowie die Umsetzung eines kompakten integrierten Hardware-Software-Systems. Damit bildet das entwickelte Produkt eine belastbare Grundlage für weiterführende Entwicklungen und eine gezielte Anpassung an unterschiedliche Anwendungsszenarien.


Im Verlauf der Entwicklung wurden mehrere grundlegende Entscheidungen getroffen, die den Charakter, die Architektur und die technische Umsetzung von Pablo maßgeblich geprägt haben.
Im Folgenden werden diese Entscheidungen jeweils mit den betrachteten Alternativen und der jeweiligen Begründung dargestellt.

\subsection{Physischer Companion statt reiner Softwareassistent}

Eine erste zentrale Entscheidung betraf die Frage, ob Pablo als reine Softwarelösung oder als physisches Produkt umgesetzt werden soll.
Softwarebasierte Assistenten lassen sich mit geringerem Aufwand entwickeln und auf bestehenden Endgeräten betreiben.
Demgegenüber bietet ein physisches System eine deutlich höhere Präsenz im Alltag der Nutzenden.
Forschungsarbeiten im Bereich der sozialen Robotik belegen, dass verkörperte Systeme eine stärkere emotionale Bindung erzeugen und als vertrauenswürdiger wahrgenommen werden als rein virtuelle Agenten.
Die Entscheidung fiel daher bewusst zugunsten eines physischen Companion-Bots, um Pablo einen eigenständigen Produktcharakter zu verleihen und eine persönliche Bindung zu ermöglichen, die über die reine Funktionalität hinausgeht.

\subsection{Modular erweiterbare Plugin-Struktur}

Für die Softwarearchitektur wurde zwischen einer monolithischen Implementierung mit fest integrierten Funktionen und einer modularen Plugin-Struktur mit dynamischer Erweiterbarkeit abgewogen.
Eine monolithische Lösung wäre in der initialen Entwicklung einfacher umsetzbar gewesen, hätte jedoch jede funktionale Erweiterung an eine Veränderung des Kernprozesses gebunden.
Die Entscheidung fiel auf eine modulare API-basierte Struktur, da diese es ermöglicht, neue Funktionen nachträglich zu ergänzen, ohne bestehenden Code verändern zu müssen.
Darüber hinaus erhöht die modulare Architektur die Anpassbarkeit an unterschiedliche Zielgruppen und Nutzungsszenarien und schafft eine Grundlage für ein langfristiges Open-Source-Ökosystem.

\subsection{Raspberry Pi als zentrale Recheneinheit}

Für die zentrale Recheneinheit wurden verschiedene Plattformen in Betracht gezogen, darunter leistungsfähigere Einplatinencomputer sowie Cloud-basierte Ansätze, bei denen die gesamte Verarbeitung auf einem externen Server erfolgt.
Die Entscheidung fiel auf einen Raspberry Pi, da dieser eine gute Verfügbarkeit, eine breite Community-Unterstützung und eine ausreichende Schnittstellenvielfalt für die Anbindung von Mikrofon, Display, Lautsprecher und Mikrocontrollern bietet.
Obwohl die Rechenleistung im Vergleich zu leistungsstärkeren Alternativen begrenzt ist, reicht sie für die lokale Koordination der einzelnen Systemkomponenten aus.
Zudem eignet sich der Raspberry Pi aufgrund seiner kompakten Bauform besonders gut für den Einsatz in einem Prototyp mit begrenztem Gehäusevolumen.

\subsection{Wakeword- und Sprachsteuerung}

Für die Aktivierung und Steuerung des Systems wurden neben der Sprachsteuerung auch physische Eingabemethoden wie Tasten oder Touchscreen-Bedienung erwogen.
Die Entscheidung fiel auf eine wakeword-basierte Sprachsteuerung, da diese eine niedrigschwellige, freihändige Interaktion ermöglicht.
Nutzende können Pablo aktivieren und Befehle erteilen, ohne ihre aktuelle Tätigkeit unterbrechen zu müssen.
Dieser Ansatz unterstützt eine möglichst natürliche Zugangsmöglichkeit zum System und entspricht dem Companion-Konzept, bei dem die Interaktion intuitiv und beiläufig erfolgen soll.

\subsection{Visuelle Persona}

Bei der Gestaltung der Benutzeroberfläche wurde zwischen einer rein funktionalen Darstellung, etwa einer textuellen Statusanzeige, und einer visuellen Persona in Form eines animierten Charakters abgewogen.
Die Entscheidung fiel auf die Darstellung eines Pixel-Art-Faultiers als visuelle Identität von Pablo.
Diese Gestaltung erzeugt einen hohen Wiedererkennungswert, ermöglicht eine freundlichere und zugänglichere Interaktion und grenzt Pablo bewusst von rein technischen Assistenzsystemen ab.
Durch zustandsabhängige Animationen kann zudem der aktuelle Systemstatus auf eine intuitive und ansprechende Weise kommuniziert werden.

\subsection{Kompaktes integriertes Display}

Alternativ zur Integration eines eigenen Displays wäre eine Ausgabe ausschließlich über Sprache oder über ein externes Gerät wie ein Smartphone möglich gewesen.
Die Entscheidung für ein kompaktes, direkt im Gehäuse verbautes Display wurde getroffen, um Pablo ein lokales visuelles Feedback zu ermöglichen.
Das Display dient der Darstellung des Avatars, der Anzeige von Statusinformationen und der visuellen Begleitung der Sprachinteraktion.
Durch die Integration des Displays in das Gehäuse wird Pablo als eigenständiges, in sich geschlossenes Produkt wahrnehmbar, das keine zusätzlichen externen Geräte für die Nutzung benötigt.